\documentclass[letterpaper]{article}

\usepackage[letterpaper,margin=1in]{geometry}

% Compile from the repository root so that this path resolves to
% the pinned WSP submodule.
\input{wsp/testing/test-case-library.tex}

\begin{document}

\def\TCID{%
  TC-0006
}

\def\TCTitle{%
  TreeView rules
}

\def\TCPurpose{%
  Verify the rules of a tree: the relationships of Add, where a node is
  placed, sorting, paths, visible rows, hit testing, keyboard navigation,
  SingleSel and label-edit decisions.
}

\def\TCPriority{%
  \TCPriorityHigh
}

\def\TCRequirementRef{%
  REQ-0013, REQ-0027, REQ-0028, REQ-0029, REQ-0030
}

\def\TCDesignRef{%
  Not applicable; the rules are in the named header.
}

\def\TCFeatureRef{%
  \texttt{include/wfc/mscomctl/tree.hpp} and \texttt{tests/tree\_test.cpp} and
  the runner \texttt{tests/tc-0006-tree.ps1}
}

\def\TCTestTechnique{%
  Equivalence partitioning and boundary value analysis of the rule functions
  (ISO/IEC/IEEE 29119-4 specification-based techniques).
}

\def\TCPreconditions{%
  The library is built standalone (\texttt{cmake -S . -B out/build}, then
  \texttt{cmake --build out/build --config Debug}) with
  \texttt{WFC\_BUILD\_TESTS} on; CTest runs the test as \texttt{TC-0006-tree}.
}

\def\TCEnvironment{%
  Any target the library builds for (Windows x86, x64 and ARM64 with MSVC,
  Linux through Winelib); local runs use the x64 Debug configuration of the
  standalone build. The test needs no window, registry or network.
}

\def\TCAssumptions{%
  The rule header under test is the one built into the library; the test is
  compiled from the same source revision.
}

\def\TCInputData{%
  Relationships -1 to 5 and 2.5; a tree of three roots; four rows with depth,
  label width and state; key names (\texttt{down}, \texttt{up}, \texttt{home},
  \texttt{end}, \texttt{left}, \texttt{right}, \texttt{space}, letters).
}

\def\TCInitialState{%
  A fresh process; the test keeps no state between runs and holds no files or
  handles.
}

\def\TCProcedure{%
  \begin{enumerate}
    \item Run \texttt{mscomctl\_tree\_test} (CTest \texttt{TC-0006-tree}).
    \item Check the relationship numbers, then place nodes for each
      relationship with and without a relative.
    \item Compare texts, join paths with a separator, count the rows that fit
      in a height.
    \item Hit-test positions on, left of, above and below rows, with an
      indentation.
    \item Ask for the effect of each key on a list of rows, with and without
      check boxes, on a collapsed node and on an empty tree.
    \item Ask the SingleSel and label-edit questions.
  \end{enumerate}
}

\def\TCExpectedResult{%
  \begin{enumerate}
    \item Relationships 0 to 4 are valid (2.5 rounds to 2), others give 35607.
    \item A node without a relative is a new root; child goes after the last
      child; first, last, next and previous go where their names say.
    \item Sorting ignores case and a prefix comes first; a path joins the
      texts with the separator and skips an empty root; the rows that fit
      follow the height.
    \item Hit testing finds the row and its indentation, and nothing outside
      the labels.
    \item Down and up wrap, home and end go to the ends, a letter finds the
      next row that starts with it, right expands or goes to the first child,
      left collapses or goes to the parent, space toggles a check box only
      with check boxes.
    \item SingleSel collapses the previous node; a second click starts an
      automatic edit; the edited text is kept unless cancelled.
    \item The program prints \texttt{tree: ok} and exits with code 0.
  \end{enumerate}
}

\def\TCPostConditions{%
  The process ends; nothing is written outside CTest's own log.
}

\TCTemplate

\end{document}
